\documentclass[11pt]{article}
\usepackage[utf8]{inputenc}
\usepackage[T1]{fontenc}
\usepackage[italian]{babel}
\usepackage{geometry}
\geometry{a4paper, margin=1in}
\usepackage{xcolor}
\usepackage{listings}
\usepackage{hyperref}

\usepackage{tcolorbox}
\tcbuselibrary{breakable}
\begin{document}

\begin{tcolorbox}[colback=blue!5!white,colframe=blue!75!black,title=\textbf{DOMANDA}, breakable]
7) Compilare ed eseguire il secondo esempio
\end{tcolorbox}

\begin{tcolorbox}[colback=green!5!white,colframe=green!75!black,title=\textbf{RISPOSTA}, breakable]
Avviando poi in background il server \texttt{./serverUDP\_inc} e successivamente lanciando il client \texttt{./clientUDP\_inc}, il comportamento osservato in esecuzione è il seguente:

\begin{itemize}
    \item Il server si mette in ascolto sulla porta 35000.
    \item Il client chiede in input un numero intero (es. ho inserito 42).
    \item Il client invia l'intero al server.
    \item Il server lo riceve, calcola \texttt{request + 1} (ottenendo 43) e lo rispedisce al client.
    \item \textbf{Il client riceve la risposta dal server e stampa a video:} \texttt{[CLIENT] Contenuto: 43}. Entrambi i programmi, arrivati alla fine, terminano la loro esecuzione.
\end{itemize}
\end{tcolorbox}

\end{document}
